\chapter{A normal current without rectifiability}\label{sec:examples-diffuse-current}
Normality controls a current and its boundary by measures, but does not force its mass to lie on a set of its own dimension. The following elementary example also makes visible the distinction between Euclidean and mass norms in \cite{Sim18}, Chapter 6, Remark 2.7, printed p.~160.

\begin{example}[A nonsimple bivector and its two norms]\label{ex:nonsimple-bivector-mass}
In oriented $\mathbb R^4$, put $\zeta=e_1\wedge e_2+e_3\wedge e_4$. Its Euclidean norm is $\sqrt2$, while its mass norm is $2$. It is not simple.
\end{example}
\begin{proof}
The two basis bivectors are orthonormal, giving the Euclidean norm. Their individual mass norms are one, so the triangle inequality gives $|\zeta|_{\mathrm{mass}}\leq2$. For the reverse inequality, let $\omega_0=dx^1\wedge dx^2+dx^3\wedge dx^4$. Define the orthogonal complex structure $Je_1=e_2$, $Je_2=-e_1$, $Je_3=e_4$, $Je_4=-e_3$. Then $\omega_0(u,v)=\langle Ju,v\rangle$, whose absolute value is at most one on orthonormal pairs. Equality holds at $(e_1,e_2)$, so $\omega_0$ has comass one. As $\langle\omega_0,\zeta\rangle=2$, duality gives the opposite mass bound. Finally $\zeta\wedge\zeta=2e_1\wedge e_2\wedge e_3\wedge e_4\ne0$, whereas the square of a simple bivector is zero.
\end{proof}

\begin{example}[A diffuse finite-mass two-current]\label{ex:diffuse-two-current}
Let $Q=[0,1]^4$ and define a two-current on $\mathbb R^4$ by
\[
 T(\omega)=\int_Q\bigl(\omega_{12}(x)+\omega_{34}(x)\bigr)\,dx.
\]
Then $\|T\|=2\mathcal L^4\llcorner Q$, its polar bivector is $\zeta/2$, and $\mathbf M(T)=2$. Its support is $Q$, and its mass measure is non-atomic.
\end{example}
\begin{proof}
The bound $|\langle\omega,\zeta\rangle|\leq2\|\omega\|_*$ shows that the formula is a current and bounds its mass by $2\mathcal L^4\llcorner Q$. For any open $W$ and any smooth scalar cutoff $0\leq g\leq1$ compactly supported in $W$, the form $g\omega_0$ has comass at most one and $T(g\omega_0)=2\int_Qg\,dx$. Taking the supremum of these cutoffs yields the reverse open-set mass bound. This identifies the measure and the polar field. Every ball about a point of $Q$ meets $Q$ in positive four-dimensional measure, while tests outside $Q$ vanish; hence the support is $Q$. Singletons have zero Lebesgue measure.
\end{proof}

\begin{proposition}[The diffuse current has a finite-mass boundary]\label{prop:diffuse-current-normal}
The current in Example \ref{ex:diffuse-two-current} is normal and satisfies $\mathbf M(\partial T)\leq8$.
\end{proposition}
\begin{proof}
For a smooth compactly supported one-form $\eta=\sum_i\eta_i\,dx^i$,
\[
 \partial T(\eta)=\int_Q(\partial_1\eta_2-\partial_2\eta_1
 +\partial_3\eta_4-\partial_4\eta_3)\,dx.
\]
Apply the fundamental theorem of calculus in the differentiated coordinate and Fubini's theorem to each of the four terms. The result is a sum of eight integrals over unit three-dimensional faces, each of a single component of $\eta$, with its appropriate sign. If the comass of $\eta$ is at most one, each component is bounded by one, so the absolute value is at most eight. Taking the supremum proves the boundary mass bound. Together with $\mathbf M(T)=2$, this proves normality. The face integrals define a one-current distributed over three-dimensional faces; they are not integration currents of those faces in dimension three.
\end{proof}

\begin{lemma}[Lower-dimensional Lipschitz images have zero ambient volume]\label{lem:lipschitz-image-ambient-null}
If $m<N$ and $f:A\subset\mathbb R^m\to\mathbb R^N$ is Lipschitz, then $\mathcal L^N(f(A))=0$, where outer measure suffices if measurability has not been specified.
\end{lemma}
\begin{proof}
First suppose $A\subset[-R,R]^m$ and let $L$ be a Lipschitz constant. Cover this cube by at most $C_R\delta^{-m}$ cubes of side $\delta$. For each cube meeting $A$, choose a point of $A$ in it. Its image and the image of the portion of $A$ in that cube lie in a ball of radius $L\sqrt m\delta$ (if $L=0$, the image is a singleton). The sum of the volumes of these covering balls is at most $C\delta^{N-m}$, tending to zero. Thus the outer measure vanishes. Intersect an unbounded domain with countably many bounded cubes to obtain the general case.
\end{proof}

\begin{proposition}[Normality does not imply rectifiability]\label{prop:diffuse-current-not-rectifiable}
The current in Example \ref{ex:diffuse-two-current} is not a rectifiable two-current, even with real multiplicities.
\end{proposition}
\begin{proof}
The mass of a rectifiable two-current is concentrated on a countably two-rectifiable set $M$. Such a set is covered, up to an $\mathcal H^2$-null set, by countably many Lipschitz images from subsets of $\mathbb R^2$. Lemma \ref{lem:lipschitz-image-ambient-null} makes all these images $\mathcal L^4$-null. The remaining $\mathcal H^2$-null set is also $\mathcal L^4$-null: covers with arbitrarily small diameters satisfy $\sum_i(\operatorname{diam}U_i)^4\leq\delta^2\sum_i(\operatorname{diam}U_i)^2$. Hence $\mathcal L^4(M)=0$. But the explicitly computed mass measure $2\mathcal L^4\llcorner Q$ assigns zero mass to $M$ and total mass two to $Q$. It cannot be concentrated on $M$. Thus the current is not rectifiable, despite its finite mass and finite boundary mass.
\end{proof}

\begin{proposition}[Rectifiable normal currents can converge to a diffuse current]\label{prop:rectifiable-real-diffuse-limit}
For each positive integer $j$, let $P_j$ be the $j^2$ centers of the equal grid squares in $[0,1]^2$. Let $S_p^{12}$ be the integration current of $[0,1]^2\times\{p\}$ with orientation $e_1\wedge e_2$, and let $S_p^{34}$ be that of $\{p\}\times[0,1]^2$ with orientation $e_3\wedge e_4$. Then
\[
 T_j=j^{-2}\sum_{p\in P_j}(S_p^{12}+S_p^{34})
\]
are compactly supported rectifiable normal two-currents with real multiplicities, $\mathbf M(T_j)=2$ and $\mathbf M(\partial T_j)=8$. They converge weakly to the nonrectifiable current $T$ of Example \ref{ex:diffuse-two-current}. Their multiplicities are not integer-valued for $j>1$.
\end{proposition}
\begin{proof}
Each current is a finite real linear combination of oriented squares. Within each family the squares are disjoint. A square in one family intersects a square in the other in at most one point, a set of zero two-dimensional measure. Their positive mass measures therefore add and give $2j^2j^{-2}=2$. Each square has boundary mass four. The grid centers are in the interior of the transverse square, so boundary edges from the two families do not meet; within one family they are also disjoint. Thus the boundary masses add to eight.

For a smooth test form $\omega$, the first-family contribution is the midpoint Riemann sum for the continuous function
\[
 p\longmapsto\int_{[0,1]^2}\omega_{12}(x_1,x_2,p_1,p_2)\,dx_1dx_2.
\]
It converges to $\int_Q\omega_{12}\,dx$. The second family gives $\int_Q\omega_{34}\,dx$. Their sum is $T(\omega)$, proving weak convergence. Almost every point on any one square lies on no other square, and there the multiplicity is $j^{-2}$. For $j>1$ it is not an integer. This example is consistent with integral-current compactness: omitting integer multiplicities changes the conclusion, even with fixed compact support and uniform mass and boundary-mass bounds.
\end{proof}

\begin{proposition}[The diffuse approximation also converges in the real flat norm]\label{prop:diffuse-current-flat-rate}
For the currents $T_j$ of Proposition \ref{prop:rectifiable-real-diffuse-limit}, the real-current flat norm of Definition \ref{def:real-current-flat-norm} satisfies
\[
 \mathcal F_{\mathbb R}(T_j-T)\leq\frac{5\sqrt2}{j}.
\]
Thus even flat convergence, fixed compact support, and uniform mass and boundary-mass bounds do not preserve real rectifiability without an additional hypothesis.
\end{proposition}
\begin{proof}
For $y\in[0,1]^2$, write $S_y^{12}$ for the corresponding oriented square, and let $p_j(y)$ be the center of its grid cell. Values on grid boundaries may be assigned arbitrarily. The displacement is at most $\delta_j=\sqrt2/(2j)$. Translating $S_y^{12}$ along the segment from $y$ to $p_j(y)$ gives the homotopy decomposition of Proposition \ref{prop:current-homotopy}:
\[
 S_{p_j(y)}^{12}-S_y^{12}=\partial R_y+B_y,
 \qquad \mathbf M(R_y)\leq\delta_j,
 \qquad \mathbf M(B_y)\leq4\delta_j.
\]
Here $R_y$ is the image of the interval product with the square, and $B_y$ is the image of the interval product with its boundary. The estimates are exactly Proposition \ref{prop:current-homotopy-estimate} with Lipschitz constant one, square mass one and boundary mass four.

Define averaged currents by integrating their test-form evaluations over $y$. These evaluations are measurable (smooth on each cell apart from its boundary), and all supports lie in the fixed cube $Q$. The uniform mass bounds imply that the integrals define finite-mass currents $R$ and $B$, with masses at most $\delta_j$ and $4\delta_j$. Their boundary identity follows by evaluating $R$ on the exterior derivative of a test form and interchanging the finite integrals. The average of $S_{p_j(y)}^{12}$ is the first-family Riemann sum, while the average of $S_y^{12}$ is the first summand of $T$. Thus this family's difference is $\partial R+B$ and has real flat norm at most $5\delta_j$. Repeat for the second family and add the decompositions. The total bound is $10\delta_j=5\sqrt2/j$. No integer-coefficient filling or integer flat-norm assertion is involved.
\end{proof}

